\section{Dualité Higgs $\leftrightarrow$ zêta}
\label{sec:dualite_higgs_zeta}
\label{sec:dualite_higgs}

La triple identification de \S\ref{sec:surdetermination_un_huitieme}
conduit à une conjecture structurelle inattendue, formulée comme verrou
K4.

\subsection{Le bifurcateur canonique commun}
\label{sec:bifurcateur}

\begin{conjecture}[K4, dualité Higgs $\leftrightarrow$ zêta]
\label{conj:K4}
Il existe un objet PT canonique $B$, le \emph{bifurcateur}, tel que~:
\begin{itemize}[nosep,leftmargin=*]
\item $m_H / v = \mathrm{scale}(B) = \sper = 1/2$ (mode breathing, masse
  du Higgs) ;
\item $\Delta_N = \mathrm{coupling}(B) = \sper^{2}/2 = 1/8$ (phase
  Maslov, secteur HP-PT) ;
\item $\lambdaH = \mathrm{coupling}(B) = \sper^{2}/2 = 1/8$ (auto-couplage
  Higgs).
\end{itemize}
\end{conjecture}

Le bifurcateur canonique est identifié comme le projecteur spinoriel
$\qplus/\qminus$ du canal $p = 2$ au point fixe $\mustar = 15$.
Formellement,
\begin{equation}
B \;=\; \Pi_+ : \Hilb \,\to\, \Hilb_{\qplus}, \qquad
\Pi_+ \;=\; \tfrac{1}{2}\big(I - i J_{\mathrm{PT}}\big),
\label{eq:bifurcateur}
\end{equation}
où $J_{\mathrm{PT}}$ est la structure complexe Kähler-Fisher de
$\mathcal{M}_m^{\mathrm{db}}$ (cf.~\S\ref{sec:K3}) et $\Hilb = \Hilb_{\qplus}
\oplus \Hilb_{\qminus}$ la décomposition chirale. L'opérateur miroir
$\Pi_- = (I + i J_{\mathrm{PT}})/2$ projette sur $\qminus$, et $B$ se
restreint au canal $p = 2$ (sous-espace propre de la partition binaire
$\Tsix$).

\subsection{Le canal $p = 2$ porte les deux observables}
\label{sec:canal_p2}

Le canal $p = 2$ joue trois rôles en PT (\cite{PT_M1})~:
\begin{itemize}[nosep,leftmargin=*]
\item opérateur info / anti-info ($\gammap{2} = 0{,}867$, le plus actif) ;
\item bifurcation $\qplus/\qminus$ (chiralité de la persistance) ;
\item partition binaire $\log_2 m$ (théorème $\Tsix$, GFT).
\end{itemize}

Le Higgs habite le canal $p = 2$ (\cite{PT_M5})~:
il est neutre sous $\mathrm{SU}(3) \times \mathrm{SU}(2) \times \mathrm{U}(1)$,
les trois jauges habitent $\{3, 5, 7\}$, donc seul $p = 2$ est admissible
pour $H$.

La phase de Maslov habite aussi le canal $p = 2$~: chaque coin de la
double barrière BK est un point où le crible bascule
$\qplus \leftrightarrow \qminus$, ce qui est l'expression
opératorielle de la partition binaire $p = 2$ ($\Tsix$).

\subsection{Dérivation explicite par la chaîne $\Tone$ + $\Ttwo$ + Haar + Weyl}
\label{sec:chaine_t1t2haarweyl}

\begin{theoreme}[Identité Higgs--Maslov via la chaîne canonique]
\label{thm:chaine}
À $\sper = 1/2$,
\begin{equation}
\boxed{\;\dfrac{1}{8} \;=\; \dfrac{\sper^{2}}{2} \;=\; \dfrac{|\lambda_2(T_{30})|}{N_{\mathrm{branches}}\,\qplus/\qminus} \;=\; \dfrac{1/4}{2} \;=\; \lambdaH \;=\; \Delta N_{\mathrm{Maslov}}.\;}
\label{eq:K4}
\end{equation}
\end{theoreme}

\begin{proof}[Esquisse, chaîne en quatre étapes]
\begin{description}[nosep,leftmargin=*]
\item[$\Tone$.] $\sper = 1/2$ (théorème, \cite{PT_M1}).
\item[$\Ttwo$.] $|\lambda_2(T_{30})| = \sper^{2} = 1/4$~: la deuxième
valeur propre de la matrice de transfert primorielle $T_{30}$ vaut
$\sper^{2}$ (gap spectral, \cite{PT_M1}).
\item[Haar.] Shift $\Reop(s_\zeta) \to \Reop(s_\zeta) + 1/2$ par le
Jacobien de la mesure de Haar multiplicative
(\S\ref{sec:haar_jacobien}).
\item[Weyl.] L'ordre symétrique de Weyl divise le compte d'états par~$2$
(symétrisation $(u, p_u) \leftrightarrow (p_u, u)$ de
\S\ref{sec:cellule_planck}), ce qui transforme $\sper^{2}$ en
$\sper^{2}/2$.
\end{description}
La composition des quatre étapes donne $|\lambda_2|/2 = \sper^{2}/2 = 1/8$.
Identifiant le secteur Higgs à $\lambdaH = \sper^{2}/2$ et le secteur HP-PT
à $\Delta_N^{\mathrm{Maslov}} = \sper^{2}/2$, on obtient \eqref{eq:K4}.
\end{proof}

La quantité $1/8$ est donc structurellement le gap spectral $\Ttwo$
divisé par le nombre de branches $\qplus/\qminus$. Les deux lectures
(Higgs et Maslov) sont deux noms pour le même nombre, généré par le même
mécanisme.

\subsection{Tests d'isolation}
\label{sec:isolation}

Aucun autre couplage PT canonique n'égale $1/8$~:

\begin{table}[ht]
\centering
\begin{tabular}{@{}lll@{}}
\toprule
\textbf{Couplage PT} & \textbf{Valeur} & \textbf{$= 1/8$ ?}\\
\midrule
$\alpha_{\mathrm{EM}}$ & $7{,}30 \cdot 10^{-3}$ & non\\
$\sin^{2}\theta_W$ & $0{,}231$ & non\\
$G_F$ (GeV$^{-2}$) & $1{,}17 \cdot 10^{-5}$ & non\\
$\lambdaH$ & $\mathbf{0{,}125}$ & \textbf{oui}\\
$\Delta N_{\mathrm{Maslov}}$ & $\mathbf{0{,}125}$ & \textbf{oui}\\
\bottomrule
\end{tabular}
\caption{Tests d'isolation : aucun couplage PT autre que $\lambdaH$ et
$\Delta N_{\mathrm{Maslov}}$ n'égale $1/8$.}
\label{tab:isolation}
\end{table}

La coïncidence numérique pure est rejetée. La hiérarchie de spin
$(\sper, \sper^{2}, \sper^{2}/2)$ est dense jusqu'à $1/8$ et se referme
algébriquement au-delà ($\sper^{3} = \sper^{2}/2$ à $\sper = 1/2$).

\subsection{Tests numériques quantitatifs}
\label{sec:tests_higgs}

À l'arbre, $m_H^{\mathrm{tree}} = v \cdot \sper = 123{,}11$ GeV (en
utilisant $v = 246{,}22$ GeV). Au premier ordre perturbatif PT,
\begin{equation}
m_H^{(1)} \;=\; v \cdot \sper \cdot (1 + C_F\,\eps) \;\simeq\; 124{,}65 \text{ GeV},
\label{eq:mH_1}
\end{equation}
puis la valeur complète PT-NLO (avec corrections d'ordre supérieur,
\cite{PT_M5}) atteint $m_H^{\mathrm{NLO}} = 125{,}287$ GeV,
où les deux coefficients sont définis intrinsèquement en PT~:
\begin{itemize}[nosep,leftmargin=*]
\item $C_F = (N_c^{2} - 1)/(2 N_c) = 4/3$ est le facteur de Casimir de la
représentation fondamentale de $\mathrm{SU}(N_c)$, avec $N_c = 3$ imposé
par PT comme seule valeur de $N_c$ pour laquelle l'équation de point
fixe $\mustar = \sum p_i$ admet une solution sur les premiers actifs
(théorème $\Tzero$, \cite{PT_M3}). Géométriquement, $C_F$
compte les transitions permises de la matrice de transfert $T_3$ du
canal $p = 3$ (7 transitions sur 9 emplacements ;
\cite{PT_M5}). Statut \textbf{[Thm]} pour la
dérivation $N_c = 3 \Rightarrow C_F = 4/3$.
\item $\eps = \alpha_s\,\sper / (2\pi)$ est le petit paramètre PT du
développement perturbatif à une boucle, qui combine le couplage fort
$\alpha_s(m_Z) \simeq 0{,}1179$ et le paramètre de bifurcation
$\sper = 1/2$ avec la normalisation standard $1/(2\pi)$ d'une boucle
gluonique (\cite{PT_M5}).
Numériquement $\eps \simeq 9{,}38 \cdot 10^{-3}$ et
$C_F\,\eps \simeq 1{,}25 \cdot 10^{-2}$. Statut \textbf{[Id]} (identité
algébrique~: formule perturbative standard avec $\alpha_s$ et $\sper$ pris
à leurs valeurs PT canoniques, sans paramètre libre).
\end{itemize}

La valeur expérimentale LHC est $125{,}25$ GeV. L'erreur du premier ordre
$m_H^{(1)}$ est $\sim 0{,}5\,\%$, et l'erreur PT-NLO complet
$m_H^{\mathrm{NLO}}$ est $0{,}0295\,\%$.

\subsection{Position finale}
\label{sec:K4_position}

K4 est établi comme \textbf{conjecture forte structurellement validée}.
La coïncidence numérique pure est rejetée. Le mécanisme commun est
identifié~: bifurcation $\qplus/\qminus$ du canal $p = 2$. La chaîne de
dérivation $\Tone \to \Ttwo \to$~Haar~$\to$~Weyl utilise uniquement des
théorèmes \textbf{[Thm]} du corpus. La promotion à théorème
inconditionnel demande la fermeture cohomologique stricte K3
(actuellement \textbf{[Thm partiel]}, \S\ref{sec:limites_ouvertures}).

